% Rà soát giả thiết và đề xuất đóng góp cho bài shape vector (Huu, An, Vinh, Tran).
% Biên dịch: XeLaTeX hoặc pdfLaTeX (Overleaf). Kiểm tra số: python check_examples.py
\documentclass[12pt,a4paper]{article}
\usepackage{iftex}
\ifPDFTeX
  \usepackage[T5]{fontenc}
  \usepackage[utf8]{inputenc}
\else
  \usepackage{fontspec}
\fi
\usepackage[vietnamese]{babel}
\usepackage[margin=2.3cm]{geometry}
\usepackage{amsmath,amssymb,amsthm}
\usepackage{booktabs,array,longtable,enumitem,tikz}
\usepackage[hidelinks]{hyperref}
\setlist{nosep,leftmargin=1.8em}
\setlength{\parskip}{5pt}
\setlength{\parindent}{0pt}

\theoremstyle{plain}
\newtheorem{kq}{Định lý}
\renewcommand{\thekq}{\Alph{kq}}
\newtheorem{md}[kq]{Mệnh đề}
\newtheorem{hq}[kq]{Hệ quả}
\newtheorem{bd}{Bổ đề}
\theoremstyle{definition}
\newtheorem{nx}[kq]{Nhận xét}
\newtheorem{vd}{Ví dụ}

\newcommand{\R}{\mathbb{R}}
\newcommand{\Lip}{\operatorname{Lip}}
\newcommand{\diam}{\operatorname{diam}}
\newcommand{\Per}{\operatorname{Per}}
\newcommand{\sd}{\mathbin{\triangle}}
\newcommand{\dm}{\,dm}

\begin{document}
\begin{center}
{\large BÁO CÁO}\\[3pt]
{\Large\bfseries Giả thiết trong các mệnh đề của bài shape vector\\ và đề xuất bổ sung}\\[5pt]
Ngày 20/10/2026
\end{center}

Báo cáo đọc lại Mục 3 đến Mục 5 của bài~\cite{huu2026}. Với mỗi mệnh đề, báo cáo ghi giả thiết được dùng ở bước nào của chứng minh và có bỏ được không. Từ đó đề xuất ba kết quả thay thế với giả thiết yếu hơn, kèm ví dụ. Số hiệu như ``MĐ~3.3'', ``BĐ~4.6'', ``VD~4.1'' chỉ mệnh đề, bổ đề, ví dụ của~\cite{huu2026}. Kết quả mới được đánh số bằng chữ cái.

\section{Vấn đề}

Bài~\cite{huu2026} định nghĩa vector hình dạng $W(x^*)=(w_1(x^*),\dots,w_N(x^*))^\top$ với
\[
w_i(x^*)=\int_\Omega f_{x^*}\,\chi_{V(x^*)\cap U_i(x^*)}\dm ,\qquad f_{x^*}(x)=f(x-x^*),
\]
rồi chứng minh $W$ nửa liên tục dưới (MĐ~3.3), liên tục (MĐ~3.5, 3.9(ii)) và Lipschitz địa phương (MĐ~3.7). Các mệnh đề này dùng nhiều giả thiết: $f$ liên tục hoặc Lipschitz, $\overline\Omega$ compact, $m(\partial U_i)=0$, cùng hai điều kiện quan sát (ĐN~3.2 và~3.6). Bài không chỉ ra giả thiết nào thật sự cần. Điều này dẫn tới ba chỗ hở.

Thứ nhất, trọng số hằng từng mảnh như $f=\chi_{B(0,\rho)}$ tính rẻ và vẫn phân biệt được hai cấu hình của VD~4.1 (xem Ví dụ~\ref{vd:step}), nhưng nằm ngoài phạm vi MĐ~3.5 và~3.7 vì không liên tục.

Thứ hai, tính cần của điều kiện quan sát liên tục chỉ được minh họa bằng Nhận xét~3.4. Trong nhận xét đó dãy $V(x_n)$ được gán tùy ý. Với $\Omega=(0,1)$, đẳng thức $V(x^*)=(0,1)$ buộc $O\cap(0,1)=\emptyset$, vì mỗi điểm vật cản $o\neq x^*$ che mọi điểm nằm sau nó. Khi đó $V(x_n)=(0,1)$ với mọi $n$. Vậy không tập vật cản nào sinh ra dãy trong Nhận xét~3.4.

Thứ ba, hằng số Lipschitz ở MĐ~3.7 bằng $N\Lip(f)\,m(U)+\|f\|_\infty\sum_i M_i$, với $M_i=2M\pi D^2/d_0(x^*)+P_i$ theo HQ~4.7. Hằng số này tăng tuyến tính theo $N$, nên cận yếu đi đúng lúc ta làm mịn phân hoạch như Mục~5 đề xuất. Mục~6 của bài cũng nêu hạn chế này. Ngoài ra $m(U)$ có thể bằng $\infty$ vì $U$ chỉ được giả thiết là tập mở chứa~$0$.

\section{Đề xuất đóng góp}

\begin{enumerate}[label=(\roman*)]
\item Bảng rà soát giả thiết (Mục~\ref{sec:audit}). Mỗi giả thiết được ghi chỗ dùng trong chứng minh và kết luận có bỏ được không. Chỗ nào không bỏ được thì có phản ví dụ.
\item Định lý~\ref{kq:A}: $W$ liên tục tại $x^*$ khi $f\in L^1(\Omega-\Omega)$, các $U_i$ đo được và $m(V(x_n)\sd V(x^*))\to 0$. Định lý này bao MĐ~3.3, 3.5 và 3.9(ii) mà không cần $f$ liên tục, $f$ Lipschitz hay $m(\partial U_i)=0$.
\item Mệnh đề~\ref{md:B}: với trọng số đều, nếu $W$ liên tục dọc $(x_n)$ ở mọi mức của một dãy phân hoạch mịn dần thì $m(V(x_n)\sd V(x^*))\to 0$. Cùng với Định lý~\ref{kq:A}, điều kiện quan sát theo độ đo là điều kiện cần và đủ.
\item Mệnh đề~\ref{md:C}: cận Lipschitz với hằng số $L_f+\|f\|_\infty\big(C_V+\sum_iP_i\big)$. Hệ số $N$ ở hai số hạng đầu biến mất, và cận áp dụng được cho trọng số hằng từng mảnh.
\item Ba ví dụ: trọng số bậc thang tách hai cấu hình của VD~4.1 với khoảng cách gấp $3{,}85$ lần trọng số $f_{\rm decay}$; một trọng số bị chặn làm $W$ mất tính Lipschitz; và so sánh hằng số Lipschitz cho phân hoạch vòng và quạt.
\end{enumerate}

\section{Thiết lập và hai công cụ của tích phân Lebesgue}

\textbf{Giả thiết (S).} $\Omega\subset\R^d$ bị chặn; $\Omega$ và $\Omega-\Omega$ đo được Lebesgue (tự động đúng khi $\Omega$ mở hoặc compact, nên đúng với bản đồ gồm các đa giác); $O\subset\R^d$; $V(x)$ định nghĩa như trong~\cite{huu2026}; $U_1,\dots,U_N\subset\R^d$ đo được và rời nhau từng đôi, $U=\bigcup_iU_i$; $f\in L^1(\Omega-\Omega)$.

$V(x)$ là phần trong của một tập con của $\Omega$, nên $V(x)$ mở, đo được và $V(x)\subset\Omega$. Đặt $\tilde f:=f$ trên $\Omega-\Omega$, $\tilde f:=0$ ngoài $\Omega-\Omega$, và $F_i:=\tilde f\chi_{U_i}$. Khi đó $F_i\in L^1(\R^d)$ và $\sum_i|F_i|\le|\tilde f|$.

\textbf{Biểu diễn (R).} Với $x\in\Omega\setminus O$,
\begin{equation}\label{eq:R}
w_i(x)=\int_{\R^d}F_i(u-x)\,\chi_{V(x)}(u)\dm(u).
\end{equation}
Thật vậy, nếu $u\in V(x)\subset\Omega$ thì $u-x\in\Omega-\Omega$, nên $f_x(u)=\tilde f(u-x)$ và $\chi_{U_i(x)}(u)=\chi_{U_i}(u-x)$. Hàm $u\mapsto F_i(u-x)$ đo được vì phép tịnh tiến biến tập đo được thành tập đo được.

\begin{bd}\label{bd:1}
Cho $g\in L^1(\R^d)$.
\begin{enumerate}[label=(\alph*)]
\item $\|g(\cdot-h)-g\|_{L^1}\to0$ khi $h\to0$.
\item Với mọi $\varepsilon>0$ có $\delta>0$ sao cho $m(A)<\delta$ kéo theo $\int_A|g|\dm<\varepsilon$.
\end{enumerate}
\end{bd}
\begin{proof}
(a) Chọn $\varphi\in C_c(\R^d)$ với $\|g-\varphi\|_{L^1}<\varepsilon/3$ \cite[Định lý 3.14]{rudin}, $\operatorname{supp}\varphi\subset B(0,\rho)$. Do độ đo Lebesgue bất biến qua tịnh tiến, $\|g(\cdot-h)-\varphi(\cdot-h)\|_{L^1}=\|g-\varphi\|_{L^1}$. Với $|h|\le1$, $\|\varphi(\cdot-h)-\varphi\|_{L^1}\le m(B(0,\rho+1))\sup|\varphi(\cdot-h)-\varphi|$, và vế phải dần về $0$ vì $\varphi$ liên tục đều. Bất đẳng thức tam giác cho kết luận.
(b) Đặt $g_k:=\min(|g|,k)$. Theo định lý hội tụ đơn điệu, có $k$ với $\int(|g|-g_k)\dm<\varepsilon/2$. Lấy $\delta:=\varepsilon/(2k)$, ta có $\int_A|g|\le\int(|g|-g_k)+k\,m(A)<\varepsilon$.
\end{proof}

\section{Rà soát giả thiết}\label{sec:audit}

Bảng~\ref{tab:audit} tóm tắt kết quả rà soát. Cột cuối dựa trên các kết quả ở Mục~\ref{sec:new}.

{\small
\begin{longtable}{@{}p{1.7cm}p{3.3cm}p{4.9cm}p{5.0cm}@{}}
\caption{Giả thiết, chỗ dùng trong chứng minh và có cần hay không.}\label{tab:audit}\\
\toprule
Kết quả & Giả thiết & Dùng ở đâu & Có cần không \\
\midrule
\endfirsthead
\toprule
Kết quả & Giả thiết & Dùng ở đâu & Có cần không \\
\midrule
\endhead
Mục 3 & $\Omega$ bị chặn & $m(\Omega)<\infty$ trong (3.5), (3.6), MĐ 3.9 & Cần. \\
Mục 3 & $\Omega$, $\Omega-\Omega$ đo được (bài không nêu) & để $\int_\Omega$ và $L^1(\Omega-\Omega)$ có nghĩa & Cần bổ sung; đúng khi $\Omega$ mở hoặc compact. \\
Mục 3 & $O\subsetneq\Omega$ & chỉ ở MĐ 4.4: $|x-v_k|\le D$ & Bỏ được nếu thay $D$ bằng $\diam(\Omega\cup O)$. Khi bỏ, VD 4.3 ($O$ rời $\Omega$) mới khớp thiết lập. \\
MĐ 3.3 & $f$ liên tục & $g_n\to g$ h.k.n. & Không cần (ĐL \ref{kq:A}). \\
MĐ 3.3 & $f\ge0$ trên $U$ & $g_n\ge0$ để dùng bổ đề Fatou; kết luận $w_i\ge0$ & Chỉ cần cho dấu. Tính liên tục (mạnh hơn nửa liên tục dưới) đúng với mọi $f\in L^1$. \\
MĐ 3.3, 3.5 & ĐN 3.2 (hội tụ h.k.n.) & $g_n\to g$; $m(A_n)\to0$ trong (3.6) & Chỉ dùng hệ quả $m(V(x_n)\sd V(x^*))\to0$; điều kiện này cần (MĐ \ref{md:B}). \\
MĐ 3.3, 3.9(ii) & $m(\partial U_i)=0$ & $\chi_{U_i(x_n)}\to\chi_{U_i(x^*)}$ h.k.n. & Không cần; chỉ cần $U_i$ đo được (BĐ \ref{bd:1}(a)). \\
MĐ 3.5, 3.9 & $\overline\Omega$ compact & $f$ bị chặn trên $\overline\Omega-\overline\Omega$ & Thừa: $\Omega$ bị chặn nên $\overline\Omega$ compact (Heine--Borel). \\
MĐ 3.5, 3.9(ii) & $f$ Lipschitz & (3.5) và $J_1$ & Không cần (ĐL \ref{kq:A}). \\
MĐ 3.7 & $f$ Lipschitz & $I_1\le\Lip(f)m(E_i(y))|y-z|$; $\|f\|_\infty<\infty$ & Thay được bằng (C3) (MĐ \ref{md:C}); không bỏ hẳn được (VD \ref{vd:G}). \\
MĐ 3.7 & $E_i$ Lipschitz theo $m(\cdot\sd\cdot)$ & $I_2$ & Cần (NX 3.8). \\
MĐ 3.9(i) & $\chi_E\in L^q$; $1<p,q<\infty$ & bất đẳng thức Hölder & $\chi_E\in L^q$ luôn đúng vì $m(E)\le m(\Omega)$; Hölder đúng cả với $(p,q)=(1,\infty),(\infty,1)$. \\
MĐ 4.3 & $O$ là hợp hữu hạn đa giác & $\Sigma(x^*)$ là hợp hữu hạn đường thẳng nên có độ đo $0$ & Đủ; chưa rõ có cần không. \\
MĐ 4.4 & $r_0=d_0(x^*)/2$ & $|y_t-v_k|\ge r_0$; $[y,z]$ không chạm $O$ & Cần cho hằng số; $C_1$ giảm được một nửa (xem dưới). \\
BĐ 4.6 & $m(\partial U)=0$ & Bước 2: bỏ tập $x-v\in\partial U$ & Cần trong chứng minh này. \\
ĐL 5.4 & $p<\infty$ & $C_c$ trù mật trong $L^p$ & Cần (phản ví dụ dưới). \\
ĐL 5.4 & $\delta_N\to0$ & $\|R_N\varphi-\varphi\|_p\le m(U)^{1/p}\omega_\varphi(\delta_N)$ & Cần: ô không co lại thì hàm không hằng trên ô đó không được xấp xỉ. \\
ĐL 5.4 & $U(x^*)\subset\Omega$; $m(\partial U_i^N)=0$ & không dùng; BĐ 5.3 chỉ cần các ô rời nhau & Bỏ được; chỉ cần $m(U)<\infty$. \\
HQ 5.5 & $m(\{g_1\ne g_2\})>0$ & $\|h\|_{L^1}>0$ & Cần và đủ. \\
\bottomrule
\end{longtable}}

Một số dòng của bảng cần giải thích thêm.

Đầu Mục 3 bài viết $V(x^*)\subset\Omega\setminus\{x^*\}$. Điểm $x^*$ thuộc tập trước khi lấy phần trong (đoạn suy biến không chạm $O$), và khi $O=\emptyset$ thì $V(x^*)=\operatorname{int}\Omega\ni x^*$. Nên viết $V(x^*)\subset\Omega$. Tích phân không đổi vì $\{x^*\}$ có độ đo $0$.

Ở MĐ 4.4, đoạn $[y,z]$ nằm trong $B(x^*,r_0)$ và $|v_k-x^*|\ge2r_0$, nên góc $\theta_k$ giữa $u_0$ và $u_1$ không vượt quá $2\arcsin(1/2)=\pi/3$. Với $\theta\in[0,\pi]$ ta có $|u_0-u_1|=2\sin(\theta/2)\ge2\theta/\pi$, do đó $\theta_k\le\frac{\pi}{2}|u_0-u_1|\le\pi|y-z|/r_0$. Bài dùng $\theta_k\le\pi|u_0-u_1|$, nên hằng số đúng là $C_1=M\pi D^2/d_0(x^*)$, bằng một nửa giá trị trong bài.

Phản ví dụ cho $p=\infty$ ở ĐL 5.4: $d=1$, $U=(0,1)$, phân hoạch dyadic, $g=\chi_{(0,1/\sqrt2)}$. Điểm $1/\sqrt2$ không bao giờ là đầu mút dyadic, nên ô chứa nó có trung bình $a\in(0,1)$ và $\|R_Ng-g\|_{L^\infty}\ge\max(a,1-a)\ge1/2$ với mọi $N$.

\section{Kết quả đề xuất}\label{sec:new}

\begin{kq}\label{kq:A}
Giả sử (S). Cho $x_n,x^*\in\Omega\setminus O$ với $x_n\to x^*$ và
\begin{equation}\label{eq:A1}
m\big(V(x_n)\sd V(x^*)\big)\to0 .
\end{equation}
Khi đó $W(x_n)\to W(x^*)$. Nếu thêm $f\ge0$ h.k.n. trên $\Omega-\Omega$ thì $w_i\ge0$.
\end{kq}
\begin{proof}
Đặt $h_n:=x_n-x^*$. Theo \eqref{eq:R},
\[
w_i(x_n)-w_i(x^*)=\underbrace{\int\big[F_i(u-x_n)-F_i(u-x^*)\big]\chi_{V(x_n)}\dm}_{J_1}
+\underbrace{\int F_i(u-x^*)\big[\chi_{V(x_n)}-\chi_{V(x^*)}\big]\dm}_{J_2}.
\]
Đổi biến $v=u-x^*$ (bất biến tịnh tiến) cho $|J_1|\le\|F_i(\cdot-h_n)-F_i\|_{L^1}$, đại lượng này dần về $0$ theo Bổ đề~\ref{bd:1}(a). Vì $|\chi_{V(x_n)}-\chi_{V(x^*)}|=\chi_{V(x_n)\sd V(x^*)}$, ta có $|J_2|\le\int_{V(x_n)\sd V(x^*)}|F_i(u-x^*)|\dm(u)$, và đại lượng này dần về $0$ theo Bổ đề~\ref{bd:1}(b) áp cho $F_i(\cdot-x^*)\in L^1$ cùng \eqref{eq:A1}. Vì $N$ hữu hạn, $W(x_n)\to W(x^*)$ theo mọi chuẩn trên $\R^N$. Kết luận về dấu suy trực tiếp từ \eqref{eq:R}.
\end{proof}

Giả thiết được dùng như sau: $f\in L^1(\Omega-\Omega)$ và $U_i$ đo được để $F_i\in L^1(\R^d)$, cần cho cả hai phần của Bổ đề~\ref{bd:1}; $x_n\in\Omega$ để \eqref{eq:R} đúng tại $x_n$; \eqref{eq:A1} chỉ dùng cho $J_2$.

ĐN 3.2 kéo theo \eqref{eq:A1}: $|\chi_{V(x_n)}-\chi_{V(x^*)}|\le\chi_\Omega$ với $m(\Omega)<\infty$, nên định lý hội tụ bị chặn cho $m(V(x_n)\sd V(x^*))\to0$. Giả thiết của MĐ 3.9(ii) ($\chi_{V(x_n)}\to\chi_{V(x^*)}$ trong $L^1(\Omega)$) chính là \eqref{eq:A1}. Vậy MĐ 3.3, 3.5 và 3.9(ii) là trường hợp riêng của Định lý~\ref{kq:A}, và kết luận của MĐ 3.3 được nâng từ nửa liên tục dưới lên liên tục.

\begin{md}\label{md:B}
Giả sử (S), $f\equiv1$, $U$ mở, bị chặn, $\Omega-\Omega\subset U$, và $\{\mathcal P^N\}$ là dãy phân hoạch mịn dần của $U$ theo ĐN 5.1. Ký hiệu $w^N_i(x):=m\big(V(x)\cap(x+U^N_i)\big)$. Nếu $x_n\to x^*$ trong $\Omega\setminus O$ và $W^N(x_n)\to W^N(x^*)$ với mọi $N$, thì \eqref{eq:A1} đúng.
\end{md}
\begin{proof}
Đặt $g_n:=\chi_{V(x_n)-x_n}$ và $g:=\chi_{V(x^*)-x^*}$, xem như phần tử của $L^2(U)$; điều này hợp lệ vì $V(x)-x\subset\Omega-\Omega\subset U$. Do bất biến tịnh tiến, $\langle g_n,\chi_{U^N_i}\rangle=w^N_i(x_n)$ và $\langle g,\chi_{U^N_i}\rangle=w^N_i(x^*)$. Vậy $\langle g_n-g,\psi\rangle\to0$ với mọi $\psi\in S:=\bigcup_N\operatorname{span}\{\chi_{U^N_i}\}$.

Chứng minh ĐL 5.4 chỉ dùng $m(U)<\infty$ và $\delta_N\to0$, nên áp được cho $U$ với $p=2$: với $\varphi\in L^2(U)$, $R_N\varphi\in S$ và $\|R_N\varphi-\varphi\|_2\to0$. Vì $\|g_n\|_2,\|g\|_2\le m(\Omega)^{1/2}$,
\[
|\langle g_n-g,\varphi\rangle|\le|\langle g_n-g,R_N\varphi\rangle|+2m(\Omega)^{1/2}\|\varphi-R_N\varphi\|_2 .
\]
Chọn $N$ để số hạng thứ hai nhỏ, rồi cho $n\to\infty$; vậy $g_n\rightharpoonup g$ yếu trong $L^2(U)$. Các ô của $\mathcal P^N$ phủ $U\supset V(x)-x$, nên $\|g_n\|_2^2=m(V(x_n))=\sum_iw^N_i(x_n)\to\sum_iw^N_i(x^*)=\|g\|_2^2$. Do đó $\|g_n-g\|_2^2=\|g_n\|_2^2+\|g\|_2^2-2\langle g_n,g\rangle\to0$.

Cuối cùng, $V(x_n)\sd V(x^*)=x_n+\big[(V(x_n)-x_n)\sd(V(x^*)-x_n)\big]$. Với $m(A\sd B)=\|\chi_A-\chi_B\|_2^2$ và bất đẳng thức tam giác cho hiệu đối xứng,
\[
m\big(V(x_n)\sd V(x^*)\big)\le\|g_n-g\|_2^2+m\big((V(x^*)-x^*)\sd(V(x^*)-x_n)\big),
\]
và số hạng cuối dần về $0$ theo Bổ đề~\ref{bd:1}(a) áp cho $\chi_{V(x^*)}$.
\end{proof}

Định lý~\ref{kq:A} và Mệnh đề~\ref{md:B} cho thấy: với trọng số đều và mọi mức phân hoạch, $W$ liên tục dọc $(x_n)$ khi và chỉ khi $m(V(x_n)\sd V(x^*))\to0$. Không thể làm yếu điều kiện quan sát thêm nữa mà vẫn giữ tính liên tục ở mọi mức phân hoạch.

\begin{md}\label{md:C}
Giả sử (S), $B:=B(x^*,r)\subset\Omega\setminus O$, và
\begin{enumerate}[label=(C\arabic*)]
\item $m(V(y)\sd V(z))\le C_V|y-z|$ với mọi $y,z\in B$;
\item $m(U_i\sd(U_i+h))\le P_i|h|$ với mọi $h\in\R^d$, $i=1,\dots,N$;
\item $f\in L^\infty(\Omega-\Omega)$ và $\int_\Omega|f(u-y)-f(u-z)|\dm(u)\le L_f|y-z|$ với mọi $y,z\in B$.
\end{enumerate}
Khi đó, với mọi $y,z\in B$,
\[
\|W(y)-W(z)\|_1\le\Big(L_f+\|f\|_{L^\infty(\Omega-\Omega)}\big(C_V+\textstyle\sum_{i=1}^NP_i\big)\Big)|y-z| .
\]
\end{md}
\begin{proof}
Viết $f_y:=f(\cdot-y)$, $\chi^y_i:=\chi_{U_i}(\cdot-y)$, $\chi_y:=\chi_{V(y)}$. Với $u\in\Omega$ ta có $u-y,u-z\in\Omega-\Omega$, nên $|f_z(u)|\le\|f\|_\infty$ với hầu hết $u\in\Omega$. Từ đồng nhất thức
\[
f_y\chi^y_i\chi_y-f_z\chi^z_i\chi_z=(f_y-f_z)\chi^y_i\chi_y+f_z(\chi^y_i-\chi^z_i)\chi_y+f_z\chi^z_i(\chi_y-\chi_z)
\]
ta được $w_i(y)-w_i(z)=A_i+B_i+C_i$. Các $U_i$ rời nhau nên $\sum_i\chi^y_i\le1$ và $\sum_i\chi^z_i\le1$. Do đó $\sum_i|A_i|\le\int_\Omega|f_y-f_z|\dm\le L_f|y-z|$ theo (C3). Bất biến tịnh tiến và (C2) cho $|B_i|\le\|f\|_\infty m\big((U_i+y)\sd(U_i+z)\big)=\|f\|_\infty m\big(U_i\sd(U_i+z-y)\big)\le\|f\|_\infty P_i|y-z|$. Cuối cùng $\sum_i|C_i|\le\|f\|_\infty m(V(y)\sd V(z))\le\|f\|_\infty C_V|y-z|$ theo (C1).
\end{proof}

Các điều kiện (C1) đến (C3) được kiểm tra như sau. (C1) là MĐ 4.4 với $C_V=C_1$ khi $r\le d_0(x^*)/2$. (C2) là BĐ 4.6 với $P_i$ là chu vi của $U_i$; với phân hoạch vòng và quạt ở Mục 4.2, ô thuộc vòng $k$ có chu vi $\frac{2\pi}{N_\theta}(R_{k-1}+R_k)+\frac{2R}{N_r}$, nên $\sum_iP_i=2R(\pi N_r+N_\theta)$. (C3) đúng với $L_f=\Lip(f)\,m(\Omega)$ khi $f$ Lipschitz trên $\Omega-\Omega$. (C3) cũng đúng khi $f=\sum_kc_k\chi_{A_k}$ với mỗi $A_k$ thỏa giả thiết BĐ 4.6, vì $\int_\Omega|f(u-y)-f(u-z)|\dm\le\sum_k|c_k|\,m\big(A_k\sd(A_k+z-y)\big)$, nên ta lấy $L_f=\sum_k|c_k|\Per(A_k)$.

So với MĐ 3.7 và HQ 4.7, hằng số cũ là $N\Lip(f)\,m(U)+\|f\|_\infty\big(NC_1+\sum_iP_i\big)$, hằng số mới là $\Lip(f)\,m(\Omega)+\|f\|_\infty\big(C_1+\sum_iP_i\big)$. Phần còn phụ thuộc phân hoạch là $\sum_iP_i$, tăng theo $N_r+N_\theta$ chứ không theo $N=N_rN_\theta$.

\begin{nx}\label{nx:D}
Nếu $f$ chỉ bị chặn, ta có thể chuyển yêu cầu chính quy sang các tập. Giả sử $E_i(\cdot)=V(\cdot)\cap U_i(\cdot)$ Lipschitz theo hiệu đối xứng với hằng số $M_i$ trên $B$, và
\begin{enumerate}[label=(C4)]
\item $m\big(E_i(z)\sd(E_i(z)+h)\big)\le Q_i|h|$ với mọi $z\in B$, $h\in\R^d$.
\end{enumerate}
Đổi biến $v=u-y$ cho $w_i(y)=\int f(v)\chi_{E_i(y)-y}(v)\dm(v)$ với $v\in\Omega-\Omega$. Vì $(E_i(y)-y)\sd(E_i(z)-z)=\big[E_i(y)\sd(E_i(z)+y-z)\big]-y$,
\[
|w_i(y)-w_i(z)|\le\|f\|_\infty\Big(m\big(E_i(y)\sd E_i(z)\big)+m\big(E_i(z)\sd(E_i(z)+y-z)\big)\Big)\le\|f\|_\infty(M_i+Q_i)|y-z|.
\]
Khi $\Omega$ và $O$ là hợp hữu hạn đa giác, ta dự đoán biên của $E_i(z)$ gồm hữu hạn đoạn thẳng và cung tròn với tổng độ dài bị chặn đều theo $z\in B$; khi đó BĐ 4.6 cho (C4). Báo cáo chưa chứng minh điều này, nên (C4) vẫn để là giả thiết. Ví dụ~\ref{vd:G} cho thấy khi cả (C3) và (C4) đều sai thì kết luận có thể sai.
\end{nx}

\begin{nx}\label{nx:E}
Nếu $f=c_i$ h.k.n. trên $U_i\cap(\Omega-\Omega)$ với mọi $i$, thì theo \eqref{eq:R}, $w_i(x;f)=c_i\,w_i(x;1)$. Khi mọi $c_i\ne0$, $W(\cdot;f)=\operatorname{diag}(c_1,\dots,c_N)\,W(\cdot;1)$, nên $W(\cdot;f)$ tách đúng những cặp cấu hình mà trọng số đều tách được. Vì vậy trọng số chỉ có ích khi nó thay đổi bên trong ít nhất một ô. Với trọng số bậc thang, các bước nhảy phải nằm bên trong ô, không nằm trên biên ô.
\end{nx}

\begin{hq}
Cho $g_1,\dots,g_K\in L^1(U(x^*))$ với $m(\{g_k\ne g_l\})>0$ khi $k\ne l$. Khi đó có $N_0$ sao cho các vector $W^N(g_1),\dots,W^N(g_K)$ đôi một khác nhau với mọi $N\ge N_0$.
\end{hq}
\begin{proof}
Áp HQ 5.5 cho từng cặp trong $K(K-1)/2$ cặp và lấy $N_0$ lớn nhất.
\end{proof}
Hệ quả này phù hợp với xu hướng ở Bảng 2 của~\cite{huu2026}, nơi số ứng viên giảm về $1$ ở mức mịn nhất với $9$ đối tượng tham chiếu. Hệ quả chỉ cho sự tồn tại của $N_0$, không cho giá trị của nó.

\section{Ví dụ minh họa}

\begin{vd}[Trọng số bậc thang trong bối cảnh VD 4.1]\label{vd:step}
Giữ nguyên dữ liệu của VD 4.1: $x^*=0$, $N_r=3$, $N_\theta=8$, $\theta_0=\pi/8$, $\rho_1=3R/4$, $\rho_2=7R/8$, $\alpha_1=\pi/15$, $\alpha_2=28\pi/225$. Lấy
\[
f_{\rm step}:=\chi_{B(0,\rho_s)},\qquad \rho_s=\tfrac{13R}{16}\in(\rho_1,\rho_2)\subset(\tfrac{2R}{3},R).
\]
Như trong~\cite{huu2026}, $V(x^*)\cap U_i(x^*)$ không đổi với $i\ne17$, và $V(x^*)\cap U_{17}=U_{17}\setminus O_\ell$ sai khác một tập độ đo $0$. Vậy $w_{17}(f;O_\ell)=\int_{U_{17}}f\dm-\int_{O_\ell}f\dm$. Với $f$ chỉ phụ thuộc $r=|x|$, tọa độ cực cho $\int_{O_\ell}f\dm=\alpha_\ell\int_{\rho_\ell}^Rf(r)\,r\,dr$. Ta được:

\begin{center}
\begin{tabular}{@{}lccc@{}}
\toprule
$f$ & $\int_{O_1}f\dm$ & $\int_{O_2}f\dm$ & $|w_{17}(f;O_1)-w_{17}(f;O_2)|$ \\
\midrule
$1$ & $7\pi R^2/480$ & $7\pi R^2/480$ & $0$ \\
$f_{\rm decay}$ & $\pi R^2/576$ & $77\pi R^2/86400$ & $73\pi R^2/86400\approx0{,}00084\,\pi R^2$ \\
$f_{\rm step}$ & $5\pi R^2/1536$ & $0$ & $5\pi R^2/1536\approx0{,}00326\,\pi R^2$ \\
\bottomrule
\end{tabular}
\end{center}

Cụ thể $w_{17}(f_{\rm step};O_1)=437\pi R^2/18432$ và $w_{17}(f_{\rm step};O_2)=497\pi R^2/18432$. Trong ví dụ này, khoảng cách giữa hai vector với $f_{\rm step}$ gấp $3{,}85$ lần khoảng cách với $f_{\rm decay}$. Trọng số $f_{\rm step}$ không liên tục nên MĐ 3.5 và 3.7 không áp được. Định lý~\ref{kq:A} áp được vì $f_{\rm step}\in L^1$, và Mệnh đề~\ref{md:C} áp được với $L_f=\Per(B(0,\rho_s))=13\pi R/8$, $\|f\|_\infty=1$. Bước nhảy $r=\rho_s$ nằm bên trong $U_{17}$, đúng như Nhận xét~\ref{nx:E} yêu cầu. Tính $w_i(x^*;f_{\rm step})$ chỉ cần diện tích của $V(x^*)\cap U_i(x^*)\cap B(x^*,\rho_s)$, cùng loại phép tính với trọng số đều. Các con số trên được kiểm tra bằng \texttt{check\_examples.py}.
\end{vd}

\begin{figure}[ht]
\centering
\newcommand{\panel}[3]{%
\begin{tikzpicture}[scale=0.85]
  \fill[gray!15] (0:2) arc (0:45:2) -- (45:3) arc (45:0:3) -- cycle;
  \foreach \r in {1,2,3} \draw[gray!60,dashed] (0,0) circle (\r);
  \foreach \a in {0,45,...,315} \draw[gray!60,dashed] (0,0) -- (\a:3);
  \fill[black!70] (#1:#3) arc (#1:#2:#3) -- (#2:3) arc (#2:#1:3) -- cycle;
  \draw[blue,thick] (0:2.4375) arc (0:45:2.4375);
  \fill (0,0) circle (1.6pt) node[below left] {$x^*$};
  \node at (22.5:3.45) {$U_{17}$};
\end{tikzpicture}}
\panel{16.5}{28.5}{2.25}\qquad\panel{11.3}{33.7}{2.625}

(a) $O_1$\hspace{6cm}(b) $O_2$
\caption{Hai cấu hình của VD 4.1. Ô $U_{17}$ tô xám nhạt, vật cản tô đậm. Cung xanh là bước nhảy $r=13R/16$ của $f_{\rm step}$. $O_1$ cắt đĩa $B(0,13R/16)$, còn $O_2$ nằm ngoài đĩa, nên $f_{\rm step}$ tách hai cấu hình dù diện tích bằng nhau.}
\end{figure}

\begin{vd}[Trọng số bị chặn làm $W$ mất tính Lipschitz]\label{vd:G}
Lấy $d=1$, $G:=\bigcup_{k\ge1}(2^{-k},2^{-k}+4^{-k})\subset(0,3/4)$, $\Omega:=(-1,1)\cup(2+G)$, $O:=\emptyset$, $x^*:=0$, $U=U_1:=(-4,4)$ ($N=1$), $f:=\chi_{2+G}$. Khi đó $V(y)=\Omega$ với mọi $y$. Với $|y|<1$ ta có $U(y)\supset\Omega$, nên $E_1(y)=\Omega$ không đổi và ĐN 3.6 đúng. Ngoài ra $m(\partial U)=0$, và $f$ bị chặn, thuộc $L^1$. Mọi giả thiết của MĐ 3.7 đều đúng, trừ tính Lipschitz của $f$.

Với $|y|<1/2$, tập $2+G+y$ rời $(-1,1)$, nên $w_1(y)=m\big(\Omega\cap(2+G+y)\big)=m\big(G\cap(G+y)\big)$. Lấy $y=4^{-K}$ với $K\ge2$. Khe hở bên trái $2^{-k}$ của $G$ là $(2^{-k-1}+4^{-k-1},2^{-k})$, có độ dài $4^{-k-1}(2^{k+1}-1)$, giảm theo $k$ và lớn hơn $y$ khi $k\le K$. Vậy với $k\le K$, khoảng $(2^{-k},2^{-k}+y)$ nằm trong $G\setminus(G+y)$. Do đó
\[
w_1(0)-w_1(y)=m\big(G\setminus(G+y)\big)\ge Ky=y\log_4(1/y),
\]
nên $|w_1(y)-w_1(0)|/|y|\to\infty$ khi $K\to\infty$, và $W$ không Lipschitz tại $0$. Theo Định lý~\ref{kq:A}, $W$ vẫn liên tục tại $0$. Ở đây (C3) và (C4) đều sai, vì cả hai đại lượng đều không nhỏ hơn $m(G\setminus(G+y))\ge Ky$. Kiểm tra số cho $K=2,\dots,5$ nằm trong \texttt{check\_examples.py}.
\end{vd}

\begin{vd}[Hằng số Lipschitz cho phân hoạch vòng và quạt]\label{vd:const}
Xét trọng số đều $f\equiv1$ (descriptor của Ouyang và cộng sự) với $N_r=n$, $N_\theta=n2^k$ như Mục 4.2. Khi đó $L_f=0$ và:

\begin{center}
\begin{tabular}{@{}ccccc@{}}
\toprule
$(n,k)$ & $N$ & $\sum_iP_i/R$ & Hệ số của $C_1$ (MĐ 3.7) & Hệ số của $C_1$ (MĐ \ref{md:C}) \\
\midrule
$(3,3)$ & $72$  & $66{,}8$  & $72$  & $1$ \\
$(4,4)$ & $256$ & $153{,}1$ & $256$ & $1$ \\
$(5,5)$ & $800$ & $351{,}4$ & $800$ & $1$ \\
\bottomrule
\end{tabular}
\end{center}

Khi không có vật cản ($O=\emptyset$, như VD 4.2) thì $C_1=0$ và hai cận trùng nhau. Khi có vùng cấm bay như VD 4.3, $C_1=2M\pi D^2/d_0(x^*)$ lớn khi $x^*$ gần vật cản, và ở mức $(5,5)$ cận cũ nhân số hạng này với $800$.
\end{vd}

\section{Việc còn mở}

Nếu chứng minh được chu vi của $V(z)$ bị chặn đều theo $z$ gần $x^*$ trong trường hợp đa giác, Nhận xét~\ref{nx:D} trở thành vô điều kiện và mọi trọng số bị chặn đều cho $W$ Lipschitz địa phương. Mệnh đề~\ref{md:B} hiện cần $U\supset\Omega-\Omega$; bản địa phương với $U=B(0,R)$ nhỏ, như trong ứng dụng, chưa được làm.

\begin{thebibliography}{9}
\bibitem{huu2026} P.~T.~Huu, P.~T.~An, D.~X.~Vinh, A.-D.~Tran, Lebesgue integral-based shape vector descriptors with applications to GPS-denied UAV localization, bản thảo (2026).
\bibitem{rudin} W.~Rudin, Real and Complex Analysis, 3rd Edition, McGraw-Hill, New York, 1987.
\end{thebibliography}

\end{document}
